\\maketitle
\\begin{abstract}
    本文针对2025年CUMCM A题智能电网负荷预测与多能互补调度优化问题，建立了"Transformer预测-NSGA-II调度-强化学习优化"的三层协同框架。首先对历史负荷数据、新能源发电数据和电价数据进行预处理；其次构建了基于Self-Attention机制的Transformer负荷预测模型，并引入STL分解提升预测精度；最后建立了考虑调度成本、新能源消纳率和用户满意度的多目标优化模型，采用NSGA-II与Q-Learning相结合的求解策略。
    
    \\textbf{针对问题一}，我们建立了STL-Transformer组合预测模型。首先利用STL方法将原始负荷序列分解为趋势项$T_t$、季节项$S_t$和残差项$R_t$：
    $$F_t = T_t + S_t + R_t$$
    趋势项和季节项采用Transformer模型进行预测，残差项采用ARIMA进行修正。实验表明，该组合模型MAPE降至3.2\\%，优于单一Transformer模型（MAPE=4.1\\%）和传统LSTM模型（MAPE=5.3\\%）。
    
    \\textbf{针对问题二}，我们设计了多能互补调度环境并建立了强化学习调度模型。状态空间定义为$S_t = [P_{load}^t, P_{wind}^t, P_{solar}^t, E_{bat}^t, t_{block}]^T$，动作空间为各机组出力调整向量$a_t = [u_{gas}^t, u_{coal}^t, u_{bat}^t]^T$。奖励函数设计为：
    $$R_t = -\\alpha \\cdot C_{cost}(t) - \\beta \\cdot C_{curtail}(t) + \\gamma \\cdot S_{satisfaction}(t)$$
    经过500轮训练，Q-Learning Agent学习到了有效的调度策略。
    
    \\textbf{针对问题三}，我们建立了多目标优化调度模型。以最小化调度成本$Z_1$、最大化新能源消纳率$Z_2$和最大化用户满意度$Z_3$为目标，约束条件包括功率平衡、机组出力限幅、储能容量和爬坡速率约束。采用NSGA-II算法求解，种群大小100，迭代200代，得到32个非支配解。选取横切点方案，调度成本降低18.7\\%，新能源消纳率提升至93.4\\%。
    
    本文提出的三层协同框架为新型电力系统调度提供了有效的技术支撑。
    
    \\keywords{智能电网\quad Transformer\quad Q-Learning\quad NSGA-II\quad 多能互补\quad 负荷预测}
\\end{abstract}
